\documentclass[11pt,reqno]{amsart}
\usepackage[T1]{fontenc}
\usepackage{amsmath,amssymb,amsthm,mathtools,booktabs}
\usepackage[margin=1in]{geometry}
\setlength{\emergencystretch}{2em}
\usepackage[colorlinks=true,linkcolor=blue,citecolor=blue,urlcolor=blue]{hyperref}
\newtheorem{theorem}{Theorem}[section]
\newtheorem{proposition}[theorem]{Proposition}
\newtheorem{lemma}[theorem]{Lemma}
\newtheorem{corollary}[theorem]{Corollary}
\theoremstyle{definition}
\newtheorem{definition}[theorem]{Definition}
\newtheorem{example}[theorem]{Example}
\newtheorem{exercise}[theorem]{Exercise}
\theoremstyle{remark}
\newtheorem{remark}[theorem]{Remark}
\newcommand{\R}{\mathbb R}
\newcommand{\F}{\mathcal F}
\newcommand{\NN}{\mathcal N}
\newcommand{\Id}{I}
\DeclareMathOperator{\diag}{diag}
\DeclareMathOperator{\arsinh}{arsinh}
\title[Finite Spin to the Harmonic Oscillator]{From Finite Spin to the Harmonic Oscillator:\\Krawtchouk Transforms, Mehler Kernels, and Convolution}
\subjclass[2020]{33C45, 22E70, 47D06, 60J60, 42A38}
\keywords{Krawtchouk polynomials, Hermite functions, harmonic oscillator, Mehler kernel, $\mathfrak{sl}_2$, contraction, Ornstein--Uhlenbeck process, Hamming scheme}
\date{Lecture notes: the symmetric-tensor and multinomial viewpoint}
\begin{document}
\begin{abstract}
The induced-matrix construction of symmetric tensor powers turns
independent finite-state systems into multinomial occupation models. For two states this gives spin, the binomial
distribution, and Krawtchouk polynomials. Recording input and output
occupations gives a joint multinomial table; the finite bilinear kernel
is its joint law divided by the product of its margins. A hypergeometric bilinear summation gives a closed finite formula;
coefficient asymptotics of that formula give Mehler's kernel. Gaussian
conjugation identifies the limiting Markov evolution with the harmonic
oscillator, and unitary continuation recovers the Fourier transform.
Cumulative Bernoulli trials specify a one-site transition within this
general framework, rather than a separate multinomial construction.
We develop the framework together with spin contraction, squeezing,
positive convolution, and rare-event limits. The final lecture derives
a system of $q-1$ oscillators from a reversible $q$-state one-site
process, with explicit normal coordinates and multivariate Krawtchouk
polynomials. Exercises are included throughout.
\end{abstract}
\maketitle
\tableofcontents
\section*{Introduction}
The organizing construction is: apply the same map independently at
$N$ sites, and then forget their labels. Linear algebra calls this a
symmetric tensor power, and probability calls its state variables
multinomial occupations. The finite multinomial foundation is the
induced-matrix framework of SFA II \cite{FeinsilverSFAII2023}; our
Gaussian and oscillator discussion continues that construction in a
fluctuation limit. In two states it gives spin-$N/2$ and the
Krawtchouk system \cite{Krawtchouk1929,FeinsilverKocik2005}.
Its central bilinear-kernel limit is obtained below by coefficient
asymptotics; it gives a Gaussian system, with Hermite polynomials,
Ornstein--Uhlenbeck transitions, and oscillator wavefunctions.
The joint input--output occupation table is the bridge: it is
multinomial before conditioning, hypergeometric under suitable fixed
margins, and Gaussian after central scaling. Distinguishing these
operations resolves several apparent ambiguities about convolution.

The route through the lectures is
\[
\begin{array}{c}
\text{SFA II: induced matrices, symmetric powers, occupation tables}\\
\downarrow\\
\text{spin, Krawtchouk polynomials, and a finite density ratio}\\
\downarrow\\
\text{closed bilinear formula, coefficient asymptotics, Mehler transition}\\
\downarrow\\
\text{oscillator, Fourier transform, squeezing, and contraction}\\
\downarrow\\
\text{convolution, degenerate kernels, and systems of oscillators}.
\end{array}
\]
The final multivariate construction is self-contained; references
\cite{Griffiths1971,Griffiths2016} provide the broader theory of
multivariate Krawtchouk polynomials and composition processes.

\subsection*{Standing conventions and limits}
$D=d/dx$, and $H=\frac12(x^2-D^2)$ is the self-adjoint oscillator on
$L^2(\R,dx)$. Differential identities are initially on Schwartz space;
exponentials are interpreted spectrally. We use
\begin{equation}\label{eq:F}
(\F f)(\xi)=\frac1{\sqrt{2\pi}}\int_\R e^{-ix\xi}f(x)\,dx,\qquad
(\F^*f)(\xi)=\frac1{\sqrt{2\pi}}\int_\R e^{+ix\xi}f(x)\,dx.
\end{equation}
Weak convergence, local lattice-density convergence, fixed-degree
polynomial convergence, and operator convergence are distinct. A
compact-window kernel calculation does not by itself imply convergence
of all finite-matrix eigenvalues. Forward references connect the finite bilinear-formula limit to the
subsequent oscillator algebra.

\section{Induced matrices, symmetric powers, and multinomial chains}\label{L:tensors}
\subsection{Independent sites and occupation vectors}
Start with $V=\mathbb C^q$ and its standard orthonormal basis
$e_1,\ldots,e_q$. Permuting the factors of $V^{\otimes N}$ leaves
invariant the space $\operatorname{Sym}^N(V)$. Its normalized occupation
basis is
\[
|\mathbf n\rangle=\binom{N}{n_1,\ldots,n_q}^{-1/2}
\sum_{\text{strings with counts }\mathbf n}
 e_{a_1}\otimes\cdots\otimes e_{a_N},\qquad\sum_a n_a=N.
\]
A one-site map $T$ induces $T^{\otimes N}$ and preserves this space.
Thus ``apply the same map independently, then forget the labels'' is
exactly the symmetric-power construction. There are
$\binom{N+q-1}{q-1}$ occupation states.
If sites are sampled independently from probabilities $p_a>0$, their
occupation law is
\begin{equation}\label{eq:multweight}
\Pi_N(\mathbf n)=\frac{N!}{\prod_an_a!}\prod_ap_a^{n_a}.
\end{equation}
In Hilbert-space language the coefficients of
$(\sum_a\sqrt{p_a}e_a)^{\otimes N}$ in this basis are
$\sqrt{\Pi_N(\mathbf n)}$. This square root of an invariant density
will become the Gaussian envelope of an oscillator wavefunction.

Let $P$ be a one-site stochastic matrix with stationary law $\mathbf p$.
The induced occupation transition $P^{[N]}$ is specified by
\begin{equation}\label{eq:occGF}
\sum_{\mathbf m}P^{[N]}(\mathbf n,\mathbf m)\mathbf z^{\mathbf m}
=\prod_a\left(\sum_bP_{ab}z_b\right)^{n_a}.
\end{equation}
Probability coordinates and orthonormal occupation coordinates should
not be silently identified. If $P$ is reversible, the symmetrized
one-site operator is $\diag(\sqrt{\mathbf p})P\diag(\sqrt{\mathbf p})^{-1}$.
Its symmetric power, in the normalized occupation basis, is
$\diag(\sqrt{\Pi_N})P^{[N]}\diag(\sqrt{\Pi_N})^{-1}$.

\subsection{Induced matrices: the SFA II framework}\label{sec:sfa-framework}
The multinomial extension is the induced-matrix construction studied
in \cite{FeinsilverSFAII2023}. We use row-stochastic matrices and
backward operators on functions. More generally, for any $q\times q$
matrix $A$, define its degree-$N$ induced matrix $A^{[N]}$ by
\[
\prod_a\left(\sum_bA_{ab}z_b\right)^{n_a}
=\sum_{|\mathbf m|=N}A^{[N]}(\mathbf n,\mathbf m)\mathbf z^{\mathbf m}.
\]
Substitution of linear forms proves the functorial identity
\begin{equation}\label{eq:induced-composition}
(AB)^{[N]}=A^{[N]}B^{[N]},\qquad I^{[N]}=I.
\end{equation}
In particular, stochastic matrices induce stochastic matrices.
An explicit form of \eqref{eq:occGF} is
\begin{equation}\label{eq:induced-table-sum}
P^{[N]}(\mathbf n,\mathbf m)
=\sum_{\substack{\sum_b c_{ab}=n_a\\\sum_a c_{ab}=m_b}}
\prod_a\frac{n_a!}{\prod_b c_{ab}!}\prod_{a,b}P_{ab}^{c_{ab}}.
\end{equation}
Each row of the table is the allocation of the sites initially in one
state. Thus the induced matrix, the independent-site occupation chain,
and the table transition are exactly the same finite object.

If $P\mathbf 1=\mathbf 1$ has a complete eigenbasis with eigenvalues
$1,r_1,\ldots,r_{q-1}$, symmetrized tensor products give eigenvalues
\begin{equation}\label{eq:induced-spectrum}
\prod_{j=1}^{q-1}r_j^{\nu_j},\qquad
|\boldsymbol\nu|\le N.
\end{equation}
The missing $N-|\boldsymbol\nu|$ factors belong to the constant mode.
A stationary one-site law $\mathbf p$ induces the stationary
multinomial law \eqref{eq:multweight}. If $P$ is reversible, its
one-site modes can be chosen orthonormal in $L^2(\mathbf p)$;
occupation eigenfunctions are then multivariate Krawtchouk polynomials.
Neither reversibility nor an orthogonal eigenbasis follows from
stochasticity alone.

For $P_t=e^{tQ}$, \eqref{eq:induced-composition} gives an occupation
semigroup with generator
\begin{equation}\label{eq:induced-generator}
(\mathcal L_N f)(\mathbf n)=
\sum_{a\ne b}n_aQ_{ab}
\bigl[f(\mathbf n-e_a+e_b)-f(\mathbf n)\bigr].
\end{equation}
The rate is the number of available sites times their one-site jump
rate. Simultaneous independent updating over a finite time and
single-site jumps in this continuous-time generator are consistent;
a discrete step that changes exactly one site is a different kernel.

\begin{remark}[Three different products]
Induced matrix composition is the exact identity
\eqref{eq:induced-composition}. Products of eigenfunctions and their
linearization coefficients are a different algebraic operation.
Positive radial or hypergroup convolution, discussed later, is a third
construction. They are related by spectral transforms in the relevant
examples, but are not interchangeable definitions of convolution.
\end{remark}

\subsection{The joint input--output table}
For stationary input and output, one site gives a pair $(A,B)$ with
probability $p_aP_{ab}$. Let $C_{ab}$ count the occurrences of each pair
among $N$ independent sites. The $q^2$ entries form a multinomial table,
whose row and column margins are the input and output occupations:
$n_a=\sum_bC_{ab}$ and $m_b=\sum_aC_{ab}$. Hence
\begin{equation}\label{eq:tableGF}
\mathbb E[\mathbf u^{\mathbf n}\mathbf v^{\mathbf m}]
=\left(\sum_{a,b}p_aP_{ab}u_av_b\right)^N,
\end{equation}
and
\begin{equation}\label{eq:tableRatio}
\frac{\Pr(\mathbf n,\mathbf m)}{\Pi_N(\mathbf n)\Pi_N(\mathbf m)}
=\frac{P^{[N]}(\mathbf n,\mathbf m)}{\Pi_N(\mathbf m)}.
\end{equation}
The density ratio in \eqref{eq:tableRatio} is the finite precursor of
the Mehler kernel. Fixing margins of a table, by contrast, leads to
hypergeometric distributions and the positive convolution considered
later. We will distinguish the unconditioned table, its conditional
laws, and its Gaussian fluctuation limit throughout.

The binary specialization, including its sampling interpretation and
normalization, is developed next. Its finite bilinear formula is then
evaluated analytically to obtain the Mehler limit.

\section{Cumulative trials and the finite probabilistic basis}\label{L:cumulative}
\subsection*{Cumulative Bernoulli trials: save, then retry}
We first specify what one update of the process means. There are $N$
labelled sites, each either a success (state $1$) or a failure (state
$0$). Suppose the input contains $k$ successes. Choose parameters
$0\le\alpha<1$ and $0<\beta<1$.
\begin{enumerate}
\item \emph{Save successes.} Independently save each of the $k$ current
successes with probability $\alpha$. A failure cannot be saved.
If $J$ successes are saved, then $J\sim\operatorname{Bin}(k,\alpha)$.
\item \emph{Retry all other sites.} Leave the $J$ saved successes
unchanged. Each of the other $N-J$ sites receives a fresh Bernoulli
trial with success probability $\beta$, independently of the saving
stage and of all other fresh trials. This includes both old failures
and old successes that were not saved.
\item \emph{Count the output.} If $R$ fresh trials succeed, the new
success count is $L=J+R$, where
$R\mid J=j\sim\operatorname{Bin}(N-j,\beta)$.
\end{enumerate}
Thus saved successes and newly obtained successes accumulate within
one update. The count need not increase: an old success that is not
saved can fail its fresh trial. On the next update the same rule is
applied anew; no success is permanently protected.

Writing $B_{m,p}(h)=\binom mhp^h(1-p)^{m-h}$, with value zero outside
$0\le h\le m$, the rule gives the explicit positive transition sum
\[
\Pr(L=\ell\mid k\text{ input successes})
=\sum_{j=0}^{\min(k,\ell)}B_{k,\alpha}(j)B_{N-j,\beta}(\ell-j).
\]
This is the saved-and-recouped-successes interpretation used here,
following the cumulative-trial viewpoint of \cite{HoareRahman1984}.
The independent one-site description below is the same update, not a
second probabilistic model.

\subsection{One-site transition and occupations}
For an input failure, the output is a success with probability
$\beta$. For an input success, it is a success either by being saved,
or by not being saved and then succeeding on the fresh trial.
The one-site matrix, with rows denoting input states and columns output
states in the order $0,1$, is therefore
\[
P_{\alpha,\beta}=
\begin{pmatrix}
1-\beta&\beta\\
(1-\alpha)(1-\beta)&\alpha+(1-\alpha)\beta
\end{pmatrix}.
\]
To obtain the symmetric kernel needed for the limit, fix $0<r<1$ and
choose
\[
\beta=\frac{1-r}{2},\qquad \alpha=\frac{2r}{1+r}.
\]
Then $\alpha(1-\beta)=r$, and the matrix reduces to
\begin{equation}\label{eq:P}
P_r=\begin{pmatrix}a&b\\b&a\end{pmatrix},\qquad
a=\frac{1+r}{2},\quad b=\frac{1-r}{2}.
\end{equation}
Here ``preserves the uniform law'' means that a site with input
probabilities $(1/2,1/2)$ has the same output probabilities:
\[
(\tfrac12,\tfrac12)P_r=(\tfrac12,\tfrac12).
\]
For functions on the two states, the matrix acts by conditional
expectation, $(P_rf)(s)=\sum_{u=0}^1P_r(s,u)f(u)$.
The function $\varepsilon(s)=(-1)^s$ has values $1,-1$, hence mean zero
under the uniform law; this is the meaning of ``centered.'' Directly,
\[
P_r\begin{pmatrix}1\\-1\end{pmatrix}
=\begin{pmatrix}a-b\\b-a\end{pmatrix}
=r\begin{pmatrix}1\\-1\end{pmatrix}.
\]
Thus $\mathbb E[\varepsilon(S_{\mathrm{out}})\mid S_{\mathrm{in}}=s]
=r\varepsilon(s)$: the constant mode has eigenvalue $1$, and the
mean-zero mode has eigenvalue $r$.

Apply $P_r$ independently to $N$ sites. Given $k$ input ones, let
$T_{r,N}(k,\ell)$ be the probability of $\ell$ output ones. Forgetting
site labels is the degree-$N$ induced-matrix construction. It gives
\begin{equation}\label{eq:transitionGF}
\sum_{\ell=0}^NT_{r,N}(k,\ell)z^\ell
=(a+bz)^{N-k}(b+az)^k.
\end{equation}
The stationary occupation law is the binomial law $\pi_N$. We define
$m_{r,N}$ to be the transition density \emph{relative to this law},
rather than the transition probability itself:
\begin{equation}\label{eq:density}
\pi_N(k)=2^{-N}\binom Nk,\qquad
m_{r,N}(k,\ell)=\frac{T_{r,N}(k,\ell)}{\pi_N(\ell)}.
\end{equation}
Thus $T_{r,N}(k,\ell)=m_{r,N}(k,\ell)\pi_N(\ell)$ and
$\sum_\ell m_{r,N}(k,\ell)\pi_N(\ell)=1$. With stationary input, it is
also the joint occupation law divided by the product of its margins:
\[
m_{r,N}(k,\ell)=
\frac{\Pr(K_{\mathrm{in}}=k,K_{\mathrm{out}}=\ell)}
{\pi_N(k)\pi_N(\ell)}.
\]
With stationary input, the probabilities of site pairs $00,01,10,11$
are $a/2,b/2,b/2,a/2$. Thus the joint occupation generating function is
\begin{equation}\label{eq:joint}
\sum_{k,\ell}\pi_N(k)T_{r,N}(k,\ell)u^kv^\ell
=\left[\frac{(1+u)(1+v)+r(1-u)(1-v)}4\right]^N.
\end{equation}
This is all the probabilistic structure required below.

\subsection{Kravchuk polynomials and the exact bilinear sum}
Define the ordinary-coefficient polynomials by
\begin{equation}\label{eq:Kgen}
(1+z)^{N-k}(1-z)^k=\sum_{n=0}^NK_n(k;N)z^n.
\end{equation}
For independent uniform input bits, the product of two such generating
functions has expectation $(1+zs)^N$. Comparing coefficients proves
\begin{equation}\label{eq:Korth}
\sum_k\pi_N(k)K_n(k;N)K_j(k;N)=\binom Nn\delta_{nj}.
\end{equation}
Write $p_{n,N}(k)=K_n(k;N)/\sqrt{\binom Nn}$ for the orthonormal
polynomials. There is also the elementary transform
\[
\sum_k\pi_N(k)u^kK_n(k;N)
=\binom Nn\left(\frac{1+u}{2}\right)^{N-n}
\left(\frac{1-u}{2}\right)^n.
\]
Apply it in both variables to
$\sum_nr^nK_n(k;N)K_n(\ell;N)/\binom Nn$.
The binomial theorem gives exactly \eqref{eq:joint}.
Recall explicitly that the kernel on the left below is
\[
m_{r,N}(k,\ell):=\frac{T_{r,N}(k,\ell)}{2^{-N}\binom N\ell},
\]
where $T_{r,N}(k,\ell)$ is the output-count probability given $k$ input
successes. Equality of coefficients therefore proves the bilinear
identity
\begin{equation}\label{eq:finmehler}
\boxed{\displaystyle
m_{r,N}(k,\ell)=\sum_{n=0}^N
\frac{r^n}{\binom Nn}K_n(k;N)K_n(\ell;N).}
\end{equation}
The model identifies the normalization and positivity; it does not yet
supply the scaling limit.


\begin{proposition}\label{prop:finmehler}
The finite bilinear kernel has transition probabilities
$T_{r,N}(k,\ell)=\pi_N(\ell)m_{r,N}(k,\ell)\ge0$ and the composition law
\[
\sum_j\pi_N(j)m_{r,N}(k,j)m_{s,N}(j,\ell)=m_{rs,N}(k,\ell).
\]
For $r=e^{-t}$, $T_{r,N}=e^{tG_N}$, where $G_N$ is the independent-bit
generator in Lecture~\ref{L:limits}.
\end{proposition}
\begin{proof}
A product of $n$ distinct site signs has conditional multiplier $r^n$,
so symmetrizing gives
\begin{equation}\label{eq:tensoreigen}
\sum_\ell T_{r,N}(k,\ell)K_n(\ell;N)=r^nK_n(k;N).
\end{equation}
Positivity comes from the sampling construction; composition follows
from orthogonality or from $P_rP_s=P_{rs}$. A single bit flips at rate
$1/2$, giving the asserted continuous-time generator.
\end{proof}

\section{The hypergeometric bilinear formula}\label{L:kernels}
\subsection{The general identity and its finite specialization}
For $t\notin\{0,-1,-2,\ldots\}$ the bilinear identity is
\begin{equation}\label{eq:general-bilinear}
\begin{aligned}
&\sum_{n\ge0}\frac{(t)_n}{n!}w^n
{}_2F_1(-n,X;t;v){}_2F_1(-n,Y;t;v)\\
&\quad=(1-w)^{X+Y-t}A^{-X-Y}
{}_2F_1\left(X,Y;t;\frac{wv^2}{A^2}\right),\qquad
A=1-(1-v)w.
\end{aligned}
\end{equation}
It is first an identity near $w=0$, with the branches taking value $1$
there; Appendix~\ref{app:bilinear} gives a short proof.

To put $t=-N$, first set $X=-k$, $Y=-\ell$. Both terminating
polynomials then involve denominator factors only up to
$\min(k,\ell)\le N$. Passage to $t=-N$ is legitimate: in the series on
the left all $n>N$ terms vanish, while the right-hand series terminates
before a denominator zero. More explicitly, each factor on the left
is a polynomial in $n$ of bounded degree with coefficients analytic
near $t=-N$, so convergence is locally uniform for $|w|<1$.
Since $(-N)_n/n!=(-1)^n\binom Nn$, we obtain
\begin{equation}\label{eq:finite-v}
\begin{aligned}
&\sum_{n=0}^N(-1)^n\binom Nn w^n
{}_2F_1(-n,-k;-N;v){}_2F_1(-n,-\ell;-N;v)\\
&\quad=(1-w)^{N-k-\ell}A^{k+\ell}
{}_2F_1\left(-k,-\ell;-N;\frac{wv^2}{A^2}\right).
\end{aligned}
\end{equation}
If $k+\ell>N$, apparent denominators cancel. The right side is the
polynomial specified by the finite left side, including at removable
singularities.

\subsection{The closed finite kernel: the starting point for the limit}
In our convention,
\[
K_n(k;N)=\binom Nn{}_2F_1(-n,-k;-N;2).
\]
Take $v=2$, $w=-r$ in \eqref{eq:finite-v}. Equations
\eqref{eq:finmehler} and \eqref{eq:finite-v} give
\begin{equation}\label{eq:closed-kernel}
\boxed{\begin{aligned}
m_{r,N}(k,\ell)
={}&(1+r)^{N-k-\ell}(1-r)^{k+\ell}\\
&\times{}_2F_1\left(-k,-\ell;-N;
-\frac{4r}{(1-r)^2}\right).
\end{aligned}}
\end{equation}
This closed bilinear expression is our asymptotic tool. Its parameters
$-k,-\ell,-N$ all grow, so fixed-parameter hypergeometric limits do not
apply. We evaluate it through its exact coefficient representation:
\begin{equation}\label{eq:coefficient}
m_{r,N}(k,\ell)
=\frac{[z^\ell](a+bz)^{N-k}(b+az)^k}
{2^{-N}\binom N\ell}.
\end{equation}
There is a direct algebraic conversion, independent of the model:
for nonzero $c$, the binomial theorem gives
\[
[z^\ell](1+cz)^{N-k}(1+dz)^k
=\binom N\ell c^\ell
{}_2F_1(-k,-\ell;-N;1-d/c).
\]
Set $c=b/a$, $d=a/b$ and restore the factor $a^{N-k}b^k$.
Then $1-d/c=-4r/(1-r)^2$, and division by
$2^{-N}\binom N\ell$ gives exactly \eqref{eq:closed-kernel}.
Thus \eqref{eq:coefficient} is an algebraic representation of the
closed bilinear formula. Its coefficient asymptotics below evaluate
that formula without assuming a joint central limit theorem.

\section{From the closed finite formula to the Mehler kernel}\label{L:jointlimit}
The central coordinate and its mesh size are
\begin{equation}\label{eq:coordinates}
x_{k,N}=\frac{N-2k}{\sqrt{2N}},\qquad
\delta_N=\sqrt{\frac2N}.
\end{equation}
The limiting probability measure is
$d\mu(x)=g(x)dx$, where $g(x)=\pi^{-1/2}e^{-x^2}$; its variance is $1/2$.


\begin{theorem}[Kravchuk bilinear kernel to Mehler]\label{thm:bilinear-limit}
Fix $0<r<1$. For any sequences $k_N,\ell_N$ such that
$x_{k_N,N}\to x$ and $x_{\ell_N,N}\to y$,
\begin{equation}\label{eq:directmehler}
\boxed{\displaystyle
m_{r,N}(k_N,\ell_N)\longrightarrow
m_r(x,y)=\frac1{\sqrt{1-r^2}}
\exp\left(\frac{2rxy-r^2(x^2+y^2)}{1-r^2}\right).}
\end{equation}
Convergence is uniform on compact central windows, and locally uniform
in $r\in(0,1)$.
\end{theorem}
\begin{proof}
We evaluate the numerator and denominator of \eqref{eq:coefficient}.
All error estimates below are uniform when the scaled coordinates
belong to a fixed compact set.

\smallskip\noindent\emph{Step 1: the stationary denominator.}
Stirling's formula in the central window gives
\begin{equation}\label{eq:binomial-asym}
2^{-N}\binom N\ell
=\sqrt{\frac2{\pi N}}\exp(-y_N^2)(1+o(1)),
\qquad y_N=x_{\ell,N}.
\end{equation}
The prefactor is $\delta_Ng(y_N)$, as required by the lattice spacing.

\smallskip\noindent\emph{Step 2: a coefficient integral.}
Put $x_N=x_{k,N}$ and write the numerator as
\begin{equation}\label{eq:integral}
C_N=\frac1{2\pi}\int_{-\pi}^{\pi}
e^{-i\ell\theta}(a+be^{i\theta})^{N-k}
(b+ae^{i\theta})^k\,d\theta.
\end{equation}
Near $\theta=0$, logarithms continuous from $0$ satisfy
\[
\begin{aligned}
\log(a+be^{i\theta})&=ib\theta-\tfrac12ab\theta^2+O(|\theta|^3),\\
\log(b+ae^{i\theta})&=ia\theta-\tfrac12ab\theta^2+O(|\theta|^3).
\end{aligned}
\]
Their combined linear coefficient before coefficient extraction is
\[
\mu_N=(N-k)b+ka=\frac N2+r\left(k-\frac N2\right),
\qquad
\mu_N-\ell=\sqrt{\frac N2}(y_N-rx_N).
\]
Consequently, for $\theta=u/\sqrt N$, the exponent in the integrand is
\begin{equation}\label{eq:local-exponent}
\frac{i(y_N-rx_N)u}{\sqrt2}
-\frac{1-r^2}{8}u^2+O\left(\frac{|u|^3}{\sqrt N}\right).
\end{equation}
This is a direct saddle-point expansion of the coefficient form of
\eqref{eq:closed-kernel}.

\smallskip\noindent\emph{Step 3: justify localization.}
Both factors have the same modulus squared:
\[
|a+be^{i\theta}|^2=|b+ae^{i\theta}|^2
=1-(1-r^2)\sin^2(\theta/2).
\]
Thus the modulus of their product in \eqref{eq:integral} is at most
\[
\exp\left[-\frac N2(1-r^2)\sin^2(\theta/2)\right].
\]
Since $\sin(|\theta|/2)\ge |\theta|/\pi$ on $[-\pi,\pi]$, the scaled
integrand is bounded by the integrable function
$\exp[-(1-r^2)u^2/(2\pi^2)]$, independently of $N,k,\ell$.
A fixed-$u$ expansion followed by dominated convergence therefore
justifies the Gaussian integral, without using a probabilistic limit
theorem. Uniformity follows by the same bound and uniform expansion
on compact sets.

\smallskip\noindent\emph{Step 4: evaluate and divide.}
Equations \eqref{eq:integral}--\eqref{eq:local-exponent} give
\begin{equation}\label{eq:numerator-asym}
C_N=\sqrt{\frac{2}{\pi N(1-r^2)}}
\exp\left[-\frac{(y_N-rx_N)^2}{1-r^2}\right](1+o(1)).
\end{equation}
Dividing by \eqref{eq:binomial-asym},
\[
m_{r,N}(k,\ell)
=\frac1{\sqrt{1-r^2}}
\exp\left[y_N^2-\frac{(y_N-rx_N)^2}{1-r^2}\right](1+o(1)).
\]
The exponent equals
$[2rx_Ny_N-r^2(x_N^2+y_N^2)]/(1-r^2)$, proving the theorem.
\end{proof}

Combining the theorem with the closed bilinear formula makes the
hypergeometric asymptotic explicit:
\begin{equation}\label{eq:hypergeo-limit}
\begin{aligned}
&(1+r)^{N-k_N-\ell_N}(1-r)^{k_N+\ell_N}
{}_2F_1\left(-k_N,-\ell_N;-N;-\frac{4r}{(1-r)^2}\right)\\
&\hspace{2em}\longrightarrow
\frac1{\sqrt{1-r^2}}
\exp\left(\frac{2rxy-r^2(x^2+y^2)}{1-r^2}\right).
\end{aligned}
\end{equation}
The coefficient representation has handled the large cancellations in
the terminating hypergeometric polynomial. Expanding its summands
individually would not justify this limit.

\subsection{The limiting transition and its Gaussian conjugation}
Put $w(x)=g(x)=\pi^{-1/2}e^{-x^2}$. The asymptotic formula just proved
gives the limiting conditional density
\[
m_r(x,y)w(y)=\frac1{\sqrt{\pi(1-r^2)}}
\exp\left[-\frac{(y-rx)^2}{1-r^2}\right].
\]
It is normalized and has the form
$Y=rx+\sqrt{1-r^2}Z$, $Z\sim\mu$. Define the joint Gaussian density
\begin{equation}\label{eq:jointgauss}
g_r(x,y)=w(x)w(y)m_r(x,y)
=\frac{\exp[-(x^2-2rxy+y^2)/(1-r^2)]}{\pi\sqrt{1-r^2}}.
\end{equation}
This is an identification of the analytic limit, not an appeal to a
joint central limit theorem to obtain it.
The corresponding operator is
\begin{equation}\label{eq:directOU}
(T_rf)(x)=\int f(rx+\sqrt{1-r^2}\,z)\,d\mu(z)
=\int\frac{g_r(x,y)}{w(x)w(y)}f(y)\,d\mu(y).
\end{equation}
One has $T_rT_s=T_{rs}$. For $r=e^{-t}$ this is the
Ornstein--Uhlenbeck semigroup with generator $G=\frac12D^2-xD$.
For negative $r$, the same formula is a reflected OU transition rather
than positive-time evolution in that semigroup.

\subsection{From a probability kernel to a wavefunction kernel}
The unitary map $U:L^2(\mu)\to L^2(dx)$,
$Uf=\pi^{-1/4}e^{-x^2/2}f$, gives the symmetric kernel
\begin{equation}\label{eq:wavefromratio}
\sqrt{w(x)}\frac{g_r(x,y)}{w(x)w(y)}\sqrt{w(y)}
=\frac1{\sqrt{\pi(1-r^2)}}
\exp\left[-\frac{1+r^2}{2(1-r^2)}(x^2+y^2)+\frac{2rxy}{1-r^2}\right].
\end{equation}
The conjugated generator is
\begin{equation}\label{eq:directconjugation}
UGU^{-1}=e^{-x^2/2}(\tfrac12D^2-xD)e^{x^2/2}
=\tfrac12D^2-\tfrac12x^2+\tfrac12=-(H-\tfrac12\Id).
\end{equation}
Thus the square roots of the limiting stationary margins provide the
Gaussian envelopes, and the kernel is $r^{H-1/2}$. Multiplying by
$\sqrt r$ gives $r^H$. The next lectures identify its Hermite
wavefunctions, Fourier phases, and spectral Mehler expansion explicitly.

\begin{exercise}
Use the coefficient integral to verify \eqref{eq:directmehler}. Check the order of the factors in
\eqref{eq:directconjugation}, including the ground-energy constant.
\end{exercise}
\section{Finite spin and Krawtchouk polynomials}\label{L:spin}

\subsection{The spin representations}
Fix $N\ge1$ and a basis $e_0,\dots,e_N$. Define
\[
R_Ne_j=(N-j)e_{j+1},\qquad L_Ne_j=je_{j-1},\qquad
\Lambda_Ne_j=(N-2j)e_j,
\]
with out-of-range terms zero. Then
\begin{equation}\label{eq:finbr}
[L_N,R_N]=\Lambda_N,\qquad [\Lambda_N,R_N]=-2R_N,\qquad
[\Lambda_N,L_N]=2L_N.
\end{equation}
This is the irreducible representation of dimension $N+1$; on
homogeneous polynomials $e_j=\xi^{N-j}\eta^j$ it is
$R=\eta\partial_\xi$, $L=\xi\partial_\eta$,
$\Lambda=\xi\partial_\xi-\eta\partial_\eta$. Comparing with
\eqref{eq:quadbr}, the dictionary
\begin{equation}\label{eq:dict}
\mathsf L\leftrightarrow-L_N,\qquad\mathsf R\leftrightarrow R_N,\qquad
\rho\leftrightarrow-\Lambda_N
\end{equation}
matches the bracket relations exactly. This is an isomorphism of
abstract Lie algebras only: the finite and quadratic representations
have different adjoint structures and are not unitarily equivalent.

With $B_N=\diag\binom Nj$, conjugation by $B_N^{1/2}$ gives orthonormal
coordinates in which
\[
\widehat R_N|j\rangle=\sqrt{(N-j)(j+1)}\,|j+1\rangle,\qquad
\widehat L_N|j\rangle=\sqrt{j(N-j+1)}\,|j-1\rangle,
\]
and $\widehat L_N=\widehat R_N^*$. Thus $J_+=\widehat L_N$,
$J_-=\widehat R_N$, $J_z=\frac12\Lambda_N$ is the spin-$\frac N2$
representation, with $m=\frac N2-j$.

The Kravchuk polynomials and their normalization have already been
established in \eqref{eq:Kgen}--\eqref{eq:Korth}. We use the same
$p_{n,N}=K_n/\sqrt{\binom Nn}$. Their duality is
$\binom NkK_n(k;N)=\binom NnK_k(n;N)$.

\subsection{The Krawtchouk transform}
Let $(\Phi_N)_{nk}=K_n(k;N)$. From \eqref{eq:Kgen} one finds
$\Phi_N^2=2^N\Id$; equivalently $\Phi_N$ is the matrix of the
substitution $(\xi,\eta)\mapsto(\xi+\eta,\xi-\eta)$ in the spin
representation \cite{FeinsilverKocik2005}. The orthogonal, symmetric
version is
\begin{equation}\label{eq:Q}
Q_N=2^{-N/2}B_N^{-1/2}\Phi_NB_N^{1/2},\qquad
(Q_N)_{kn}=\sqrt{\pi_N(k)}\,p_{n,N}(k),\qquad Q_N^2=\Id.
\end{equation}
The \emph{symmetric Krawtchouk matrix} $S_N=\Phi_NB_N$ is symmetric but
not orthogonal; its spectrum is a different and harder problem
\cite{FeinsilverFitzgerald1996}.

\begin{proposition}\label{prop:Qdiag}
Let $M_N=\frac N2\Id-\frac12(\widehat R_N+\widehat L_N)$. Then
$M_NQ_N=Q_N\diag(0,1,\dots,N)$.
\end{proposition}
\begin{proof}
This is the three-term recurrence of the $K_n$ in the variable $k$,
rewritten in orthonormal coordinates (equivalently, the difference
equation of Lecture~\ref{L:limits}).
\end{proof}
Thus $\widehat R_N+\widehat L_N=2J_x$ plays the role of a position
operator, and $M_N$ is a ``finite oscillator'' Hamiltonian with exactly
equally spaced spectrum $0,1,\dots,N$; compare
\cite{AtakishievSuslov1990,AtakishievWolf1997}.

\begin{exercise}
Check $\Phi_2^2=4\Id$ and $Q_2^2=\Id$ explicitly, and verify
Proposition~\ref{prop:Qdiag} for $N=2$.
\end{exercise}
\begin{exercise}
Show that the Casimir $J_x^2+J_y^2+J_z^2$ equals
$\frac N2(\frac N2+1)\Id$ and explain why no finite-dimensional
representation can satisfy $[a,a^\dagger]=\Id$.
\end{exercise}

%======================================================================
\section{Hermite functions, the oscillator, and the Fourier transform}\label{L:hermite}

\subsection{Two-parameter Hermite polynomials}
\begin{definition}
For $\tau\in\R$ define polynomials $H_n(x,\tau)$ by
\begin{equation}\label{eq:Hgen}
e^{zx-\tau z^2/2}=\sum_{n\ge0}H_n(x,\tau)\frac{z^n}{n!}.
\end{equation}
Write $H_n(x)=H_n(x,1)$ for the probabilists' Hermite polynomials.
\end{definition}
Replacing $z$ by $z\sqrt\tau$ shows $H_n(x,\tau)=\tau^{n/2}H_n(x/\sqrt\tau)$
for $\tau>0$; in particular $H_n(x,\tfrac12)=2^{-n/2}H_n(\sqrt2x)$.
Differentiating \eqref{eq:Hgen} with respect to $x$, to $z$, and to
$\tau$ gives
\begin{equation}\label{eq:Hrec}
\partial_xH_n=nH_{n-1},\qquad
H_{n+1}=xH_n-n\tau H_{n-1},\qquad
\partial_\tau H_n+\tfrac12\partial_x^2H_n=0.
\end{equation}
Combining the first two, $\tau H_n''-xH_n'+nH_n=0$.

\subsection{Hermite functions and ladder operators}
Put
\[
\psi_n(x)=2^nH_n(x,\tfrac12)e^{-x^2/2},\qquad
H=\tfrac12(x^2-D^2).
\]
\begin{proposition}\label{prop:ladder}
For all $n\ge0$, with $\psi_{-1}=0$,
\[
(x+D)\psi_n=2n\psi_{n-1},\qquad
(x-D)\psi_n=\psi_{n+1},\qquad
H\psi_n=(n+\tfrac12)\psi_n.
\]
Consequently $\psi_n=(x-D)^ne^{-x^2/2}$ and
\begin{equation}\label{eq:psigen}
\sum_{n\ge0}\psi_n(x)\frac{u^n}{n!}=e^{-x^2/2+2ux-u^2}.
\end{equation}
\end{proposition}
\begin{proof}
Write $H_n=H_n(x,\frac12)$. By the product rule
$(x+D)(f e^{-x^2/2})=f'e^{-x^2/2}$ and
$(x-D)(fe^{-x^2/2})=(2xf-f')e^{-x^2/2}$. The first and second
identities follow from \eqref{eq:Hrec} with $\tau=\frac12$, since
$2xH_n-nH_{n-1}=2H_{n+1}$. Then
$2H=(x+D)(x-D)-1$ gives $2H\psi_n=(2n+1)\psi_n$.
The generating function is \eqref{eq:Hgen} with $z=2u$, $\tau=\frac12$.
\end{proof}

The normalized Hermite functions
\[
h_n=\frac{\psi_n}{\pi^{1/4}\sqrt{2^nn!}}
=\frac{H_n(\sqrt2x)e^{-x^2/2}}{\pi^{1/4}\sqrt{n!}}
\]
form an orthonormal basis of $L^2(\R,dx)$. With
\[
a=\frac{x+D}{\sqrt2},\qquad a^\dagger=\frac{x-D}{\sqrt2},\qquad
\NN=a^\dagger a=H-\tfrac12\Id,
\]
one has $[a,a^\dagger]=\Id$, $ah_n=\sqrt n\,h_{n-1}$,
$a^\dagger h_n=\sqrt{n+1}\,h_{n+1}$, and $\NN h_n=nh_n$.
The self-adjoint operator $H$ is defined spectrally by this basis.

\subsection{The Fourier transform as a quarter period}
\begin{theorem}\label{thm:fourier}
$\F h_n=(-i)^nh_n$ and $\F^*h_n=i^nh_n$. Consequently
\begin{equation}\label{eq:Fosc}
\F=e^{-i\pi\NN/2},\qquad
e^{-i\pi H/2}=e^{-i\pi/4}\F,\qquad
e^{i\pi H/2}=e^{i\pi/4}\F^*.
\end{equation}
\end{theorem}
\begin{proof}
The Gaussian integral
$(2\pi)^{-1/2}\int e^{zx-x^2/2}dx=e^{z^2/2}$, valid for complex $z$,
applied to \eqref{eq:psigen} gives
\[
\F^*\Big(\sum\psi_n\tfrac{u^n}{n!}\Big)(\xi)
=e^{-u^2}e^{(i\xi+2u)^2/2}=e^{-\xi^2/2+2i\xi u+u^2}
=\sum\psi_n(\xi)\frac{(iu)^n}{n!}.
\]
Comparing coefficients yields $\F^*\psi_n=i^n\psi_n$, and taking
inverses $\F\psi_n=(-i)^n\psi_n$. Both sides of \eqref{eq:Fosc} are
unitary and agree on an orthonormal basis.
\end{proof}

\begin{remark}[Signs]
The sign in \eqref{eq:F} matters. If the positive-exponent transform is
called $\F$, then $\F$ and $\F^*$ must be exchanged everywhere, and the
formula reads $e^{i\pi H/2}=\sqrt i\,\F$ with $\sqrt i=e^{i\pi/4}$.
In particular, with \eqref{eq:F}, the correct statement is
$e^{i\pi H/2}=\sqrt i\,\F^{-1}$. A common source of error is a lost
factor $2$ in the mixed term: the exponent is $2i\xi u$, not $i\xi u$.
\end{remark}
More generally $e^{-i\theta\NN}$ is the fractional Fourier transform of
angle $\theta$. Its discrete counterpart built from Krawtchouk functions
is studied in \cite{AtakishievWolf1997}; we meet it in
Lecture~\ref{L:limits}.

\subsection{The quadratic $\mathfrak{sl}_2$}
On $\mathcal S(\R)$ put
\begin{equation}\label{eq:quad}
\mathsf R=\tfrac12x^2,\qquad \mathsf L=\tfrac12D^2,\qquad \rho=xD+\tfrac12.
\end{equation}
Using $[D,x]=1$ one checks
\begin{equation}\label{eq:quadbr}
[\mathsf L,\mathsf R]=\rho,\qquad[\mathsf L,\rho]=2\mathsf L,\qquad
[\rho,\mathsf R]=2\mathsf R,\qquad H=\mathsf R-\mathsf L.
\end{equation}
The closure of $\rho$ is skew-adjoint and generates the unitary
dilations $(e^{s\rho}f)(x)=e^{s/2}f(e^sx)$. Thus
$\{\mathsf L,\mathsf R,\rho\}$ generate rotations (Fourier),
dilations (squeezing), and, in imaginary time, heat semigroups. These
quadratic operators preserve parity, so this representation splits into
even and odd pieces; it is very different from a finite irreducible
representation, a point we return to in Lecture~\ref{L:limits}.

\begin{exercise}
Verify \eqref{eq:quadbr}, and show that
$K_+=\frac12(a^\dagger)^2$, $K_-=\frac12a^2$, $K_0=\frac12H$ satisfy
$[K_0,K_\pm]=\pm K_\pm$, $[K_-,K_+]=2K_0$. Find the lowest $K_0$
weights on even and odd functions.
\end{exercise}
\begin{exercise}
Show that $\F x\F^{-1}=-D_\xi/i$ in the sense
$\widehat{xf}=iD_\xi\widehat f$ and $\widehat{Df}=i\xi\widehat f$, and
deduce $\F(x-D)\F^{-1}=-i(x-D)$. Use this to give a second proof of
Theorem~\ref{thm:fourier}.
\end{exercise}

%======================================================================
\section{The oscillator semigroup and Mehler's formula}\label{L:mehler}

\subsection{The heat kernel of $H$}
\begin{theorem}\label{thm:heat}
For $s>0$, $e^{-sH}$ is the integral operator on $L^2(\R,dx)$ with kernel
\begin{equation}\label{eq:heat}
K_s(x,y)=\frac1{\sqrt{2\pi\sinh s}}
\exp\left[-\frac{(x^2+y^2)\cosh s-2xy}{2\sinh s}\right].
\end{equation}
\end{theorem}
\begin{proof}
The kernel in \eqref{eq:wavefromratio} has already been obtained from
the finite bilinear formula. Set $r=e^{-s}$ and multiply it by
$e^{-s/2}$; the identities
$1-e^{-2s}=2e^{-s}\sinh s$ and
$(1+e^{-2s})/(1-e^{-2s})=\coth s$ give \eqref{eq:heat}.
Proposition~\ref{prop:onbasis} verifies the action on every Hermite
basis vector. The Gaussian kernel is Hilbert--Schmidt, so agreement
on that basis identifies the bounded operator with $e^{-sH}$.
\end{proof}

\begin{proposition}\label{prop:onbasis}
For $t<0$, with $K_{-t}$ as in \eqref{eq:heat},
\[
\int_\R K_{-t}(x,y)\psi_n(y)\,dy=e^{t(n+1/2)}\psi_n(x).
\]
\end{proposition}
\begin{proof}
Integrate the kernel against the generating function \eqref{eq:psigen}.
A Gaussian integral gives
\[
e^{t/2}e^{-x^2/2+2e^tux-e^{2t}u^2}
=\sum_n\frac{u^n}{n!}e^{t(n+1/2)}\psi_n(x).
\]
Compare coefficients.
\end{proof}

\subsection{Positive time and analytic continuation}
For real $t>0$ the spectral operator $e^{tH}$ is unbounded, with domain
\[
\mathcal D(e^{tH})=\Big\{\sum c_nh_n:\ \sum e^{2t(n+1/2)}|c_n|^2<\infty\Big\},
\]
which contains every $h_n$, and $e^{tH}h_n=e^{t(n+1/2)}h_n$. However, the
kernel \eqref{eq:heat} with $s=-t<0$ does \emph{not} give a convergent
integral, even on $\psi_n$: after multiplying by $\psi_n(y)$ the
coefficient of $y^2$ in the exponent is
$\frac12(\coth t-1)=(e^{2t}-1)^{-1}>0$. For complex $t$ with
$\operatorname{Re}t<0$, the kernel
\[
\frac1{\sqrt{-2\pi\sinh t}}
\exp\left[\frac{(x^2+y^2)\cosh t-2xy}{2\sinh t}\right]
\]
is analytic, with the root continued from the positive root on
$t<0$. Its boundary value at $t=i\pi/2$ is
$(2\pi)^{-1/2}e^{i\pi/4}e^{ixy}$, recovering \eqref{eq:Fosc}.
Boundary values are understood as oscillatory integrals or strong
limits, not as absolutely convergent kernels on all of $L^2$.

\subsection{Mehler's formula}
\begin{theorem}[Mehler \cite{Mehler1866}]\label{thm:mehler}
For real $x,y$ and complex $|r|<1$,
\begin{equation}\label{eq:mehler}
\sum_{n\ge0}\frac{r^n}{n!}H_n(x)H_n(y)
=\frac1{\sqrt{1-r^2}}
\exp\left(\frac{2rxy-r^2(x^2+y^2)}{2(1-r^2)}\right),
\end{equation}
the root being the branch equal to $1$ at $r=0$. The series converges
absolutely and locally uniformly in the disk, and its radius of
convergence is exactly $1$.
\end{theorem}
\begin{proof}
The closed Gaussian expression has already been obtained in
\eqref{eq:directmehler}. In the variance-one coordinates
$\xi=\sqrt2x$, $\eta=\sqrt2y$, it is the right side of
\eqref{eq:mehler}, as a kernel relative to
$d\gamma(\eta)=(2\pi)^{-1/2}e^{-\eta^2/2}d\eta$.
Gaussian integration against $e^{u\eta-u^2/2}$ gives
$e^{ru\xi-r^2u^2/2}$. Comparing coefficients shows that $H_n$ has
multiplier $r^n$. Orthogonality and completeness identify the displayed
Hermite expansion for $0<r<1$. Cauchy estimates on the Hermite
generating function give subexponential growth of normalized modes on
compact sets, so the series converges absolutely and locally uniformly
for $|r|<1$. Analytic continuation extends the identity to this disk.
The singularities at its boundary show that the radius is $1$.
\end{proof}

\begin{corollary}\label{cor:mehlerL2}
For $0<r<1$,
\begin{equation}\label{eq:mehlernorm}
\sum_{n\ge0}r^nh_n(x)h_n(y)
=\frac1{\sqrt{\pi(1-r^2)}}
\exp\left[-\frac{1+r^2}{2(1-r^2)}(x^2+y^2)+\frac{2r}{1-r^2}xy\right],
\end{equation}
and this is the kernel of $r^{\NN}$ on $L^2(\R,dx)$. With $r=e^{-s}$,
multiplying by $e^{-s/2}$ gives \eqref{eq:heat}.
\end{corollary}
\begin{proof}
Substitute $h_n$ into \eqref{eq:mehler} at $(\sqrt2x,\sqrt2y)$.
The kernel is square-integrable, hence defines a bounded operator,
which agrees with $r^{\NN}$ on the Hermite basis.
\end{proof}

\begin{remark}[Which operator is Mehler's sum?]
The bare sum \eqref{eq:mehler} is a kernel relative to $\gamma$:
\[
(T_rf)(x)=\int M_r(x,y)f(y)\,d\gamma(y)
\]
satisfies $T_rH_n=r^nH_n$, so $T_r=r^{L}$ with $L=-D^2+xD$ on
$L^2(\gamma)$. This is the Ornstein--Uhlenbeck semigroup. The unitary
$f\mapsto(2\pi)^{-1/4}e^{-x^2/4}f(x)$ followed by the dilation
$x\mapsto\sqrt2x$ intertwines $L$ with $\NN$. Thus Mehler's sum is
$r^{H-1/2}$ \emph{after} Gaussian conjugation and rescaling.
\end{remark}

\begin{exercise}
Show $\|e^{tH}\|=e^{t/2}$ for $t<0$, and that $e^{tH}$ is trace class
with $\operatorname{tr}e^{tH}=1/(2\sinh(-t/2))$. Verify this by
integrating \eqref{eq:heat} along the diagonal.
\end{exercise}
\begin{exercise}
Show that the semigroup property $K_s*K_{s'}=K_{s+s'}$ (composition of
kernels) reduces to the addition formulas for $\cosh$ and $\sinh$.
\end{exercise}

%======================================================================
\section{Gaussian kernels as squeezed oscillator semigroups}\label{L:gauss}
Consider, for $a>0$ and real $b$, the operator $A_{a,b}$ with kernel
\[
G_{a,b}(x,y)=e^{-a(x^2+y^2)+bxy}.
\]
\begin{theorem}\label{thm:gauss}
Suppose $|b|<2a$ and set
\[
\omega=\sqrt{4a^2-b^2},\qquad r=\frac{b}{2a+\omega},\qquad
C=\sqrt{\frac{\pi(1-r^2)}{\omega}}.
\]
Let $H_\omega=\frac12(-D^2+\omega^2x^2)$,
$\NN_\omega=H_\omega/\omega-\frac12$, and
$h^{(\omega)}_n(x)=\omega^{1/4}h_n(\sqrt\omega x)$. Then
\[
G_{a,b}(x,y)=C\sum_{n\ge0}r^nh^{(\omega)}_n(x)h^{(\omega)}_n(y).
\]
Thus $A_{a,b}$ is trace class with eigenvalues $Cr^n$;
$A_{a,b}=Cr^{\NN_\omega}$ for $r>0$,
$A_{a,b}=CP|r|^{\NN_\omega}$ for $r<0$ (with $Pf(x)=f(-x)$), and
$A_{a,0}$ is rank one.
\end{theorem}
\begin{proof}
One checks $a=\frac{\omega(1+r^2)}{2(1-r^2)}$ and
$b=\frac{2\omega r}{1-r^2}$, and then applies \eqref{eq:mehlernorm} in
the variables $\sqrt\omega x,\sqrt\omega y$. For $r<0$ use
$Ph_n=(-1)^nh_n$.
\end{proof}
No translation is needed, because the quadratic form has no linear
terms. The rescaling is the unitary dilation
$V_\omega=e^{\frac12(\log\omega)\rho}$, with $\NN_\omega=V_\omega\NN V_\omega^{-1}$.
So $A_{a,b}$ is a function of a \emph{squeezed} oscillator, not of $H$
itself. For $0<b<2a$, writing $r=e^{-\tau}$,
\begin{equation}\label{eq:gaussheat}
a=\frac\omega2\coth\tau,\qquad b=\frac\omega{\sinh\tau},\qquad
A_{a,b}=\sqrt{\frac{2\pi\sinh\tau}{\omega}}\,e^{-\tau H_\omega/\omega}.
\end{equation}
In the Lie-algebraic form of \eqref{eq:quad},
$e^{-\tau H_\omega/\omega}=\exp[\tau(\omega^{-1}\mathsf L-\omega\mathsf R)]$.

\begin{example}\label{ex:sqrt8}
For $a=\frac32$, $b=1$ one finds $\omega=\sqrt8$, $r=3-\sqrt8$,
$\tau=2\log(1+\sqrt2)$. Kernels of exactly this type arise as candidate
central limits in the spectral analysis of the symmetric Krawtchouk
matrices of Lecture~\ref{L:spin}, whose spectra were studied in
\cite{FeinsilverFitzgerald1996}. The theorem identifies the Gaussian
$e^{-\sqrt8x^2/2}$ as the \emph{exact} ground state of such a limiting
operator. Transferring this to the finite matrices is a separate,
genuinely asymptotic problem: compact-window kernel asymptotics give
lower bounds via Rayleigh quotients, but convergence of the full spectrum
requires global estimates.
\end{example}

\begin{exercise}
Show that if $|b|>2a$ the kernel does not define a bounded operator on
$L^2(\R)$.
\end{exercise}

%======================================================================
\section{From spin to oscillator: contraction and diffusion}\label{L:limits}

\subsection{Contraction to the Heisenberg algebra}
Set $a_N=\widehat L_N/\sqrt N$, $a_N^\dagger=\widehat R_N/\sqrt N$,
$\NN_N=\frac12(N\Id-\Lambda_N)$. Then
\begin{equation}\label{eq:contr}
[a_N,a_N^\dagger]=\Id-\frac{2}{N}\NN_N,\qquad
[\NN_N,a_N]=-a_N,\qquad[\NN_N,a_N^\dagger]=a_N^\dagger.
\end{equation}
For fixed $j$, $a_N|j\rangle\to\sqrt j\,|j-1\rangle$ and
$a_N^\dagger|j\rangle\to\sqrt{j+1}\,|j+1\rangle$. Under $|j\rangle\mapsto h_j$
this is the oscillator representation of Lecture~\ref{L:hermite}. The
convergence holds on each fixed finite set of levels, near the
highest-weight end; it is not uniform over all levels, as the trace
obstruction to $[a,a^\dagger]=\Id$ demands.

The quadratic algebra \eqref{eq:quad} is then obtained \emph{after} the
limit, as quadratic expressions in $a,a^\dagger$. The spin algebra does
not contract to the quadratic $\mathfrak{sl}_2$; it contracts to the
Heisenberg algebra, from which $\mathfrak{sl}_2$ is rebuilt.

\subsection{Krawtchouk to Hermite}
Use the central coordinate and mesh
\begin{equation}\label{eq:scale}
x_k=\frac{N-2k}{\sqrt{2N}},\qquad\delta_N=\sqrt{2/N}.
\end{equation}
\begin{theorem}\label{thm:KH}
If $x_{k_N}\to x$, then for each fixed $n$,
\[
n!\,(2N)^{-n/2}K_n(k_N;N)\to H_n(x,\tfrac12),\qquad
p_{n,N}(k_N)\to\frac{H_n(\sqrt2x)}{\sqrt{n!}},
\]
and $\delta_N^{-1/2}(Q_N)_{k_Nn}\to h_n(x)$.
\end{theorem}
\begin{proof}
Put $z=u/\sqrt{2N}$ in \eqref{eq:Kgen}. With $N-k=\frac N2+\sqrt{N/2}\,x_k$
and $k=\frac N2-\sqrt{N/2}\,x_k$,
\[
(N-k)\log(1+z)+k\log(1-z)=xu-\tfrac{u^2}4+o(1),
\]
locally uniformly in complex $u$. Cauchy estimates give coefficientwise
convergence to \eqref{eq:Hgen} with $\tau=\frac12$. The de Moivre--Laplace
local limit $\pi_N(k_N)/\delta_N\to\pi^{-1/2}e^{-x^2}$ gives the last
statement.
\end{proof}
Limit relations of this type between discrete and continuous classical
families are discussed systematically in \cite{GodoyEtAl1998}.

\subsection{Ehrenfest to Ornstein--Uhlenbeck}
The Ehrenfest urn chain has generator
\begin{equation}\label{eq:GN}
(G_Nf)(k)=\frac{N-k}2\big(f(k+1)-f(k)\big)+\frac k2\big(f(k-1)-f(k)\big).
\end{equation}
It is reversible for $\pi_N$, and $G_Np_{n,N}=-np_{n,N}$. Equivalently,
each of $N$ independent bits flips at rate $\frac12$.
\begin{theorem}\label{thm:OU}
For smooth $F$ and $f(k)=F(x_k)$,
$G_Nf(k)\to(GF)(x)$ where $G=\frac12D^2-xD$. The invariant measures
converge to $d\mu=\pi^{-1/2}e^{-x^2}dx$. The unitary
$U:L^2(\mu)\to L^2(dx)$, $Uf=\pi^{-1/4}e^{-x^2/2}f$, satisfies
\[
-UGU^{-1}=H-\tfrac12\Id=\NN.
\]
At finite $N$, $U_Nf(k)=\sqrt{\pi_N(k)}f(k)$ gives $-U_NG_NU_N^{-1}=M_N$.
\end{theorem}
\begin{proof}
Taylor expansion with step $\delta_N$: the drift is
$\frac{N-2k}{2}\cdot(-\delta_N)\to-x$ and the diffusion coefficient is
$\frac N2\cdot\frac{\delta_N^2}{2}\to\frac12$. The conjugation is a
direct computation: $e^{-x^2/2}(\frac12D^2-xD)e^{x^2/2}=\frac12D^2-\frac12x^2+\frac12$.
\end{proof}
Note that the Hermite polynomials appear in two different stochastic
roles. In the fixed-$N$ urn above they are eigenfunctions of a
\emph{stationary} diffusion. Along a simple random walk $X_m$, by
contrast, the polynomials $K_n$ with order parameter $m$ give discrete
martingales, whose limits $H_n(B_\tau,\tau)$ are Brownian martingales by
the heat equation in \eqref{eq:Hrec}. These are distinct scalings.

\subsection{Finite fractional transforms}
Define $\mathfrak F_{N,\theta}=Q_N\diag(e^{-in\theta})Q_N^T=e^{-i\theta M_N}$.
These are unitary, and on each fixed limiting mode they converge to
$e^{-i\theta\NN}$; at $\theta=\frac\pi2$ the limit is $\F$. This is the
fractional Fourier--Krawtchouk transform of \cite{AtakishievWolf1997}.
Note that $Q_N$ itself is an involution while $\F$ has order four;
$Q_N$ is the analogue of a reflection in the spin group, not of the
Fourier transform.

\begin{exercise}
Show that with the opposite coordinate $x_k=(2k-N)/\sqrt{2N}$ every
limit in Theorem~\ref{thm:KH} acquires a factor $(-1)^n$.
\end{exercise}

%======================================================================
\section{Convolution: what survives the limit}\label{L:conv}

\subsection{Fixed margins and Krawtchouk convolution}
The table viewpoint distinguishes independent transitions from a
fixed-count flip. If a fixed set of $k$ occupied sites is intersected
with a uniformly sampled set of $\ell$ flip sites, the four occupancy
categories are conditioned on these two margins. The overlap count
is hypergeometric. Mixing over a binomial flip count removes that
conditioning and returns the independent-site transition.

Identify $\{0,1\}^N$ with subsets of $\{1,\dots,N\}$; the group law is
symmetric difference. If independent uniform subsets have sizes
$k,\ell$, the size of their sum is $k+\ell-2J$, with $J$ hypergeometric.
This defines the radial (hypergroup) convolution
\begin{equation}\label{eq:hyper}
\delta_k*_N\delta_\ell=\sum_j
\frac{\binom kj\binom{N-k}{\ell-j}}{\binom N\ell}\,\delta_{k+\ell-2j}.
\end{equation}
Its identity is $\delta_0$. The normalized characters
$\chi_{n,N}(k)=K_n(k;N)/\binom Nn$ satisfy the product formula
\begin{equation}\label{eq:char}
\int\chi_{n,N}\,d(\delta_k*_N\delta_\ell)=\chi_{n,N}(k)\chi_{n,N}(\ell),
\end{equation}
so the transform $\widehat\nu(n)=\sum_k\chi_{n,N}(k)\nu(k)$ turns $*_N$
into multiplication. These are the eigenvalues of the Hamming
association scheme \cite{Delsarte1973,MacWilliamsSloane1977}; the
product formula is a special case of the Krawtchouk addition theorem
\cite{Dunkl1976}.

\subsection{Central inputs: collapse}
\begin{proposition}\label{prop:collapse}
If $x_{k_N}\to x$ and $x_{\ell_N}\to y$, then the pushforward of
$\delta_{k_N}*_N\delta_{\ell_N}$ to the $x$-coordinate converges to $\mu$,
independently of $x,y$.
\end{proposition}
\begin{proof}
The mean of $x_M$ is $\sqrt{2/N}\,x_{k_N}x_{\ell_N}\to0$ and its variance
tends to $\frac12$; apply the hypergeometric central limit theorem.
\end{proof}
Two central inputs thus produce the invariant Gaussian, forgetting both.
Spectrally, $\chi_{n,N}(k_N)=O(N^{-n/2})$, and the factor $N^{n/2}$ needed
to obtain Hermite polynomials appears twice on the right of
\eqref{eq:char} but once on the left. The reason is geometric: the
identity $k=0$ has moved to $x_0=\sqrt{N/2}\to\infty$.

\subsection{One macroscopic input: the Mehler kernel}
\begin{theorem}\label{thm:mixed}
Let $x_{k_N}\to x$ and $\ell_N/N\to\frac{1-r}2$ with $-1<r<1$. Then
\[
x_M\Longrightarrow rx+\sqrt{1-r^2}\,Z,\qquad Z\sim\mu,
\]
and $\chi_{n,N}(\ell_N)\to r^n$ for each fixed $n$.
\end{theorem}
Now one factor in \eqref{eq:char} stays of order one, and the product
formula has a nontrivial limit: the Ornstein--Uhlenbeck kernel. There
is an exact finite version. The radial binomial measure
\[
\nu_{r,N}(\ell)=\binom N\ell\Big(\frac{1-r}2\Big)^\ell\Big(\frac{1+r}2\Big)^{N-\ell}
\]
(flip each bit with probability $\frac{1-r}{2}$) has character
multipliers $r^n$, so convolution by $\nu_{r,N}$ \emph{is} the finite
Mehler operator of Proposition~\ref{prop:finmehler}, and
$\nu_{r,N}*_N\nu_{q,N}=\nu_{rq,N}$. In the limit this becomes
\[
\int m_r(x,z)m_q(z,y)\,d\mu(z)=m_{rq}(x,y),
\]
i.e.\ the semigroup law $e^{-sH}e^{-tH}=e^{-(s+t)H}$.

\subsection{Rare flips and neighboring Poisson limits}
The macroscopic input $\ell_N/N\to(1-r)/2$ is a fixed-density scaling,
not a Poisson limit. If each site is flipped with probability $p_N$
and $Np_N\to\lambda$, the number of flips tends to
$\operatorname{Poisson}(\lambda)$, but a tight number of flips moves
the central coordinate by an amount tending to zero.
For $p_N=\lambda/N$, $N$ repeated steps have degree-$n$ multiplier
$(1-2\lambda/N)^{Nn}\to e^{-2\lambda n}$, giving $e^{2\lambda G}$
or, after Gaussian conjugation, $e^{-2\lambda\NN}$.
A different rare-event limit changes the stationary law itself:
for biased-binomial Krawtchouk polynomials one may use the generating
function $(1+(1-p)z)^k(1-pz)^{N-k}$. With $p\to0$, $Np\to\lambda$,
and $k$ fixed, it tends to $(1+z)^ke^{-\lambda z}$, a Charlier generating
function. Central Gaussian and rare-event Poisson scalings are thus
neighboring but distinct regimes of occupation models.

\subsection{Other products}
Ordinary convolution on $\R$ is addition of independent variables,
diagonalized by characters $e^{ix\xi}$. The OU operation instead
contracts one input and adds independent noise. The Hermite polynomials
satisfy the binomial identity
\[
H_n(ax+by,\tfrac12)=\sum_j\binom nja^jb^{n-j}H_j(x,\tfrac12)H_{n-j}(y,\tfrac12),
\qquad a^2+b^2=1,
\]
which splits degrees rather than multiplying characters. On polynomials
one can also define the Wick product
$H_m(\cdot,\frac12)\diamond H_n(\cdot,\frac12)=H_{m+n}(\cdot,\frac12)$,
equivalently $f\diamond g=W((W^{-1}f)(W^{-1}g))$ with $W=e^{-D^2/4}$.
It is an algebraic product, not a positive convolution.

\begin{remark}[An intertwining question]
One may ask for an intertwiner between ordinary convolution with the
Fourier transform and Krawtchouk convolution with the Krawtchouk
transform. The discussion above shows that the question needs
qualification: a finite convolution theorem exists, and a natural
Mehler/OU limit exists, but under central scaling Krawtchouk convolution
does not tend to ordinary convolution. Any proposed intertwiner must
specify the product, the measure, and the topology to be preserved.
\end{remark}

\begin{exercise}
Compute the mean and variance of $k+\ell-2J$ and prove
Proposition~\ref{prop:collapse} and Theorem~\ref{thm:mixed}.
\end{exercise}
\begin{exercise}
Verify that $\nu_{r,N}$ has character multipliers $r^n$, using
\eqref{eq:Kgen} with $z$ replaced by a suitable ratio.
\end{exercise}

%======================================================================
\section{The degenerate Gaussian kernels: $b=\pm2a$}\label{L:end}
We return to Lecture~\ref{L:gauss} at the boundary of the region
$|b|<2a$, where the oscillator picture breaks down.

\subsection{$b=2a$: free heat flow}
Here $G_{a,2a}(x,y)=e^{-a(x-y)^2}$, so $A_{a,2a}$ is ordinary convolution
with a Gaussian. By \eqref{eq:F},
$\F(A_{a,2a}f)(\xi)=\sqrt{\pi/a}\,e^{-\xi^2/(4a)}\F f(\xi)$, hence
\[
A_{a,2a}=\sqrt{\pi/a}\;e^{D^2/(4a)}.
\]
It is bounded with norm $\sqrt{\pi/a}$, but not compact; its spectrum is
$[0,\sqrt{\pi/a}]$, purely continuous. Ordinary convolution appears
exactly when the confining quadratic form degenerates.

\subsection{$b=-2a$: heat flow with reflection}
Here $G_{a,-2a}(x,y)=e^{-a(x+y)^2}$ and
\[
A_{a,-2a}=\sqrt{\pi/a}\;e^{D^2/(4a)}P,
\]
with $P$ commuting with the heat operator. It equals the previous
operator on even functions and its negative on odd functions; its
spectrum is $[-\sqrt{\pi/a},\sqrt{\pi/a}]$.

\subsection{Approach to the endpoints}
As $b\uparrow2a$ with $a$ fixed, $\omega\to0$, $r\to1$, and
$\tau/\omega\to1/(2a)$. In \eqref{eq:gaussheat} the potential
$\frac{\omega^2x^2}{2}$ disappears, the kinetic part tends to
$D^2/(4a)$, and the scalar tends to $\sqrt{\pi/a}$. The uniform Schur
bound
\[
\sup_x\int e^{-a(x^2+y^2)+bxy}\,dy=\sqrt{\frac\pi a}\,
\sup_xe^{-(4a^2-b^2)x^2/(4a)}=\sqrt{\frac\pi a}\qquad(|b|\le2a)
\]
and pointwise convergence on compactly supported functions give strong
convergence $A_{a,b}\to A_{a,2a}$ on $L^2$. Norm convergence fails,
since the interior operators are compact and the endpoint operator is
not. The case $b\downarrow-2a$ follows by composing with $P$. For
$|b|>2a$ the form is indefinite, and the kernel no longer defines a
bounded operator.

\begin{exercise}
Show that the eigenvalues $Cr^n$ of $A_{a,b}$ accumulate densely in
$[0,\sqrt{\pi/a}]$ as $b\uparrow2a$, consistent with the spectrum of
the endpoint operator.
\end{exercise}

%======================================================================

\section{From $q$ states to a system of oscillators}\label{L:multi}
\subsection{One-site dynamics and normal modes}
Let $q\ge2$, $d=q-1$, and let $Q$ be an irreducible continuous-time
Markov generator on $q$ states, reversible for probabilities $p_a>0$:
$p_aQ_{ab}=p_bQ_{ba}$. Choose a real orthonormal basis of $L^2(\mathbf p)$,
\begin{equation}\label{eq:modes}
u_0=1,\qquad Qu_j=-\lambda_ju_j\ (1\le j\le d),\qquad
\sum_ap_au_i(a)u_j(a)=\delta_{ij},\quad\lambda_j>0.
\end{equation}
Repeated eigenvalues cause no difficulty: choose an orthonormal basis
in each eigenspace. At time $t$, $P_t=e^{tQ}$ has nonconstant
multipliers $r_j(t)=e^{-\lambda_jt}$. Independent copies give the
occupation generator
\begin{equation}\label{eq:compgen}
(\mathcal G_Nf)(\mathbf n)=\sum_{a\ne b}n_aQ_{ab}
[f(\mathbf n-\mathbf e_a+\mathbf e_b)-f(\mathbf n)],
\end{equation}
reversible for $\Pi_N(\mathbf n)=\frac{N!}{\prod_an_a!}\prod_ap_a^{n_a}$.
This is a composition process. Multivariate Krawtchouk systems and their
composition processes are studied in \cite{Griffiths1971,Griffiths2016};
we develop the needed formulas directly.

\subsection{Whitening the multinomial fluctuations}
The raw vector $(\mathbf n-N\mathbf p)/\sqrt N$ has covariance
$\Sigma=\diag(\mathbf p)-\mathbf p\mathbf p^T$, of rank $d$.
Rather than invert a singular covariance matrix, set
\begin{equation}\label{eq:multiX}
X_{j,N}(\mathbf n)=\frac1{\sqrt{2N}}\sum_an_au_j(a),\qquad 1\le j\le d.
\end{equation}
These coordinates have mean zero and covariance $\delta_{ij}/2$.
Together with $\sum_an_a=N$, they determine the counts:
\begin{equation}\label{eq:invertX}
n_a=Np_a+\sqrt{2N}\,p_a\sum_{j=1}^dX_{j,N}u_j(a).
\end{equation}
The central limit theorem gives $\mathbf X_N\Rightarrow\mu_d$, where
\[
d\mu_d(\mathbf x)=\pi^{-d/2}e^{-|\mathbf x|^2}\,d\mathbf x.
\]
There are $q-1$ fluctuation modes, not $q$: total population is conserved.
For stationary one-site input and output separated by time $t$,
\[
\mathbb E[u_i(A)u_j(B)]=\langle u_i,P_tu_j\rangle_p
=r_j(t)\delta_{ij}.
\]
The joint central limit separates into $d$ independent Gaussian pairs.
Conditionally on input occupations tending to $\mathbf x$ in these
coordinates, output fluctuations tend to
\begin{equation}\label{eq:multiOU}
Y_j=r_j(t)x_j+\sqrt{1-r_j(t)^2}\,Z_j,\qquad \mathbf Z\sim\mu_d.
\end{equation}
This conditional assertion follows by grouping the independent output
variables according to their input state and applying a triangular-array
central limit theorem. Joint weak convergence alone would not justify
conditioning on an individual input lattice point.
Expanding \eqref{eq:compgen} gives the same limiting generator:
\begin{equation}\label{eq:multiG}
\mathcal G=\sum_{j=1}^d\lambda_j
\left(\tfrac12\partial_{x_j}^2-x_j\partial_{x_j}\right).
\end{equation}
The drift follows from $Qu_j=-\lambda_ju_j$. For the diffusion use
\[
\sum_{a,b}p_aQ_{ab}(u_i(b)-u_i(a))(u_j(b)-u_j(a))
=2\lambda_j\delta_{ij}.
\]
A local limit theorem gives the corresponding kernel-density limit.
The lattice cell volume in the normal coordinates cancels between the
joint law and the product of its marginal laws.

\subsection{Multivariate Krawtchouk polynomials}
For $\boldsymbol\alpha=(\alpha_1,\ldots,\alpha_d)$ write
$|\boldsymbol\alpha|=\sum_j\alpha_j$ and
$\boldsymbol\alpha!=\prod_j\alpha_j!$. Define ordinary coefficients by
\begin{equation}\label{eq:multiK}
G_{\mathbf n}(\mathbf z):=\prod_{a=1}^q
\left(1+\sum_{j=1}^dz_ju_j(a)\right)^{n_a}
=\sum_{|\boldsymbol\alpha|\le N}K_{\boldsymbol\alpha}(\mathbf n)
\mathbf z^{\boldsymbol\alpha}.
\end{equation}
Each coefficient is a symmetrized tensor with $\alpha_j$ sites in mode
$u_j$, all remaining sites in mode $u_0$. Independence and
orthonormality give
\[
\mathbb E_{\Pi_N}[G_{\mathbf n}(\mathbf z)G_{\mathbf n}(\mathbf w)]
=(1+\mathbf z\cdot\mathbf w)^N.
\]
Therefore
\begin{equation}\label{eq:multiorth}
\mathbb E_{\Pi_N}[K_{\boldsymbol\alpha}K_{\boldsymbol\beta}]
=B_{\boldsymbol\alpha,N}\delta_{\boldsymbol\alpha\boldsymbol\beta},
\qquad B_{\boldsymbol\alpha,N}=
\frac{N!}{(N-|\boldsymbol\alpha|)!\boldsymbol\alpha!}.
\end{equation}
The polynomials $p_{\boldsymbol\alpha,N}=K_{\boldsymbol\alpha}/
\sqrt{B_{\boldsymbol\alpha,N}}$ form a complete orthonormal basis:
there are $\binom{N+d}{d}$ of them, just as many as occupation states.
Tensor eigenvectors imply
\[
\mathcal G_Np_{\boldsymbol\alpha,N}
=-\Big(\sum_j\lambda_j\alpha_j\Big)p_{\boldsymbol\alpha,N}.
\]
The finite kernel relative to $\Pi_N$ is consequently
\begin{equation}\label{eq:multifinite}
m_{t,N}(\mathbf n,\mathbf m)=\sum_{|\boldsymbol\alpha|\le N}
\left(\prod_j r_j(t)^{\alpha_j}\right)
p_{\boldsymbol\alpha,N}(\mathbf n)p_{\boldsymbol\alpha,N}(\mathbf m).
\end{equation}
As in the binary case, this is the joint multinomial input--output
law divided by the product of its stationary margins.

When $\mathbf X_N(\mathbf n_N)\to\mathbf x$, substitute
$z_j=v_j/\sqrt{2N}$ in \eqref{eq:multiK}. Its logarithm tends locally
to $\mathbf x\cdot\mathbf v-|\mathbf v|^2/4$. Comparing fixed-degree
coefficients gives
\begin{equation}\label{eq:multiKH}
p_{\boldsymbol\alpha,N}(\mathbf n_N)\longrightarrow
\prod_{j=1}^d\frac{H_{\alpha_j}(\sqrt2x_j,1)}{\sqrt{\alpha_j!}}.
\end{equation}
This directly generalizes the binary Krawtchouk--Hermite limit.

\subsection{Product Mehler kernels and oscillator energies}
The density of \eqref{eq:multiOU} relative to $\mu_d$ is
\begin{equation}\label{eq:multimehler}
m_t(\mathbf x,\mathbf y)=\prod_{j=1}^d
\frac1{\sqrt{1-r_j(t)^2}}
\exp\left[\frac{2r_j(t)x_jy_j-r_j(t)^2(x_j^2+y_j^2)}{1-r_j(t)^2}\right].
\end{equation}
Conjugate by $U_df=\pi^{-d/4}e^{-|\mathbf x|^2/2}f$. Then
\begin{equation}\label{eq:multiHamiltonian}
\begin{aligned}
-U_d\mathcal GU_d^{-1}&=\sum_j\lambda_j\NN_j
=\mathcal H-\frac12\sum_j\lambda_j\Id,\\
\mathcal H&=\frac12\sum_j\lambda_j(-\partial_{x_j}^2+x_j^2).
\end{aligned}
\end{equation}
The normalized eigenfunctions and energies are
\[
h_{\boldsymbol\alpha}(\mathbf x)=\prod_jh_{\alpha_j}(x_j),\qquad
E_{\boldsymbol\alpha}=\sum_j\lambda_j(\alpha_j+\tfrac12).
\]
The kernel of $e^{-t\mathcal H}$, $t>0$, is
\begin{equation}\label{eq:multiheat}
\prod_{j=1}^d\left\{
\frac{1}{\sqrt{2\pi\sinh(\lambda_jt)}}
\exp\left[-\frac{(x_j^2+y_j^2)\cosh(\lambda_jt)-2x_jy_j}
{2\sinh(\lambda_jt)}\right]\right\}.
\end{equation}
Thus the continuum limit is genuinely a system of $q-1$ oscillators.
Under individual coordinate dilations $x_j=\sqrt{\lambda_j}\,\xi_j$,
$\mathcal H$ becomes the unit-mass form
$\frac12\sum_j(-\partial_{\xi_j}^2+\lambda_j^2\xi_j^2)$ with frequencies
$\lambda_j$. Couplings in raw occupation coordinates are removed by
whitening and the normal-mode transformation. If all $\lambda_j=\lambda$,
the total-excitation-$n$ energy is $\lambda(n+d/2)$ with multiplicity
$\binom{n+d-1}{d-1}$.

For Fourier phases, a common angle is generated by
$\exp[-i\theta\sum_j\NN_j]$. Physical time under the number Hamiltonian
$\sum_j\lambda_j\NN_j$ gives different angles $\lambda_jt$.
A simultaneous ordinary Fourier transform requires all those angles
to equal $\pi/2$ modulo $2\pi$; unequal frequencies do not in general
have a common quarter period.

\subsection{The algebraic limit: $\mathfrak{sl}_q$ to $d$ bosons}
Choose the orthonormal one-site Hilbert vectors
$v_0(a)=\sqrt{p_a}$ and $v_j(a)=\sqrt{p_a}u_j(a)$.
The symmetrized negative generator has eigenvalues
$0,\lambda_1,\ldots,\lambda_d$ in this basis. In its $N$th symmetric
power let $E_{ab}=c_a^\dagger c_b$, $0\le a,b\le d$, be the
number-conserving matrix units. They obey
\[
[E_{ab},E_{cd}]=\delta_{bc}E_{ad}-\delta_{ad}E_{cb},\qquad
\sum_{a=0}^dE_{aa}=N\Id.
\]
Trace-free combinations generate the finite $\mathfrak{sl}_q$
representation. Put
$b_{j,N}=E_{0j}/\sqrt N$ and $b_{j,N}^\dagger=E_{j0}/\sqrt N$.
With finitely many excitations outside mode $0$,
\begin{equation}\label{eq:multicontraction}
[b_{i,N},b_{j,N}^\dagger]=\delta_{ij}E_{00}/N-E_{ji}/N
\longrightarrow\delta_{ij}\Id.
\end{equation}
The creation coefficient is $\sqrt{(n_j+1)n_0/N}\to\sqrt{n_j+1}$;
the annihilation coefficient is $\sqrt{n_j(n_0+1)/N}\to\sqrt{n_j}$.
The finite symmetrized composition generator, expressed in normal-mode
occupations, is exactly $\sum_j\lambda_jE_{jj}$; its limit is
$\sum_j\lambda_jb_j^\dagger b_j$, the number Hamiltonian in
\eqref{eq:multiHamiltonian}. This identifies the probabilistic and
algebraic limits. As before, convergence is on fixed finite excitation
sectors, not uniformly over the entire finite representation.

\subsection{The multistate save-and-redistribute model}
This is an explicit reversible family within the induced-matrix
framework \cite{FeinsilverSFAII2023}, not a replacement for the general
multinomial model. Save a
site of type $a$ with probability $\alpha_a$; redraw every unsaved site
from probabilities $\beta_b$. This gives
\[
P_{ab}=\alpha_a\delta_{ab}+(1-\alpha_a)\beta_b,\qquad
\eta_a=\frac{\beta_a/(1-\alpha_a)}{\sum_c\beta_c/(1-\alpha_c)}.
\]
For $0\le\alpha_a<1$ and positive $\beta_a$, detailed balance follows
from $\eta_a(1-\alpha_a)\beta_b=\eta_b(1-\alpha_b)\beta_a$.
Thus occupation equilibrium is multinomial with probabilities
$\boldsymbol\eta$, not generally $\boldsymbol\beta$.
Diagonalizing this reversible one-site matrix supplies the mode
multipliers $r_j$. The limiting Gaussian kernel is
$\prod_jr_j^{\NN_j}$; for $0<r_j<1$ it can be interpolated as an OU
semigroup with times $-\log r_j$. This does not assert that every such
finite stochastic matrix itself embeds in a continuous-time chain.
It is the discrete-step counterpart of the generator construction
above, and makes the several-oscillator interpretation concrete.

\begin{example}[Three states and two oscillators]
Let $p_a=1/3$ and $Qf(a)=\frac13\sum_bf(b)-f(a)$. Take
\[
u_1=\sqrt{3/2}(1,-1,0),\qquad u_2=(1,1,-2)/\sqrt2.
\]
Both modes have $\lambda_1=\lambda_2=1$, and
\[
X_{1,N}=\frac{\sqrt3(n_1-n_2)}{2\sqrt N},\qquad
X_{2,N}=\frac{n_1+n_2-2n_3}{2\sqrt N}.
\]
The limit is $\mathcal H=\frac12(-\Delta+|\mathbf x|^2)$ on $\R^2$.
Its energy at total excitation $n$ is $n+1$, with multiplicity $n+1$.
The Mehler kernel factors into two binary kernels, even though the
finite occupation simplex is triangular rather than a product of two
binomial lattices.
\end{example}

\subsection{Scope and further directions}
The finite foundation is the induced multinomial chain of SFA II
\cite{FeinsilverSFAII2023}. The bilinear-kernel asymptotics, Gaussian conjugation,
and fixed-excitation oscillator contraction developed here are further
steps, not consequences of stochasticity alone.
Reversibility gives orthogonal real modes, a self-adjoint conjugated
Hamiltonian, and independent Mehler factors. Nonreversible one-site
chains can have Gaussian limits too, but their drift need not be
orthogonally diagonalizable and their conjugated generator need not be
a self-adjoint oscillator. Such an identification requires additional
assumptions. Nor does the raw multinomial central limit theorem imply
spectral convergence: here the exact tensor eigenbasis and
fixed-excitation contraction identify the limiting number operators.

\begin{exercise}
Prove \eqref{eq:invertX} from completeness of $u_0,\ldots,u_d$ and
\eqref{eq:multiorth} by coefficient comparison. Verify the two normal
coordinates in the three-state example.
\end{exercise}
\begin{exercise}
For $q=2$ show that \eqref{eq:multicontraction} gives
$[a_N,a_N^\dagger]=1-2\NN_N/N$. Explain why two occupation numbers
produce only one oscillator.
\end{exercise}
\appendix
\section{Proof of the general bilinear identity}\label{app:bilinear}
Initially take $t>X,Y>0$, $0<v<1$, and $w>0$ small. Euler's beta
integral represents ${}_2F_1(-n,X;t;v)$ as the average of $(1-vs)^n$
against the beta$(X,t-X)$ density. Use it for both factors in
\eqref{eq:general-bilinear} and sum
$\sum_n(t)_nz^n/n!=(1-z)^{-t}$. The resulting double integral has
integrand $[1-w(1-vs)(1-vu)]^{-t}$.

Integrating over $s$ first, and setting $c=1-w$, $d=wv$, $A=c+d$,
gives
\[
\frac1{B(Y,t-Y)}\int_0^1
u^{Y-1}(1-u)^{t-Y-1}
(c+du)^{X-t}[A+d(1-v)u]^{-X}\,du.
\]
Make the substitution $u=ch/(A-dh)$. The integral becomes
\[
\frac{c^{X+Y-t}A^{-X-Y}}{B(Y,t-Y)}\int_0^1
h^{Y-1}(1-h)^{t-Y-1}
\left(1-\frac{wv^2}{A^2}h\right)^{-X}dh.
\]
Euler's integral now gives \eqref{eq:general-bilinear}. Analytic
continuation proves it for general parameters wherever the indicated
functions are defined, initially as analytic germs at $w=0$.


\begin{thebibliography}{99}
\bibitem{AtakishievSuslov1990}
N.~M. Atakishiev and S.~K. Suslov, \emph{Difference analogs of the
harmonic oscillator}, Theoret. and Math. Phys. \textbf{85} (1990).
\href{https://doi.org/10.1007/BF01017247}{doi:10.1007/BF01017247}.

\bibitem{AtakishievWolf1997}
N.~M. Atakishiyev and K.~B. Wolf, \emph{Fractional Fourier--Kravchuk
transform}, J. Opt. Soc. Amer. A \textbf{14} (1997), no.~7, 1467.
\href{https://doi.org/10.1364/JOSAA.14.001467}{doi:10.1364/JOSAA.14.001467}.

\bibitem{CooperHoareRahman1977}
R.~D. Cooper, M.~R. Hoare, and M.~Rahman,
\emph{Stochastic processes and special functions: On the probabilistic
origin of some positive kernels associated with classical orthogonal
polynomials}, J. Math. Anal. Appl. \textbf{61} (1977), no.~1, 262--291.
\href{https://doi.org/10.1016/0022-247X(77)90160-3}{doi:10.1016/0022-247X(77)90160-3}.

\bibitem{HoareRahman1984}
M.~R. Hoare and M.~Rahman, \emph{Cumulative Bernoulli trials and
Krawtchouk processes}, Stochastic Process. Appl. \textbf{16} (1984),
no.~2, 113--139.
\href{https://doi.org/10.1016/0304-4149(84)90014-0}{doi:10.1016/0304-4149(84)90014-0}.
The publisher dates the issue February 1984; its copyright notice is 1983.

\bibitem{Delsarte1973}
P.~Delsarte, \emph{An algebraic approach to the association schemes of
coding theory}, Philips Research Reports Supplements, No.~10 (1973).

\bibitem{Dunkl1976}
C.~F. Dunkl, \emph{A Krawtchouk polynomial addition theorem and wreath
products of symmetric groups}, Indiana Univ. Math. J. \textbf{25}
(1976), 335--358.

\bibitem{Eagleson1969}
G.~K. Eagleson, \emph{A characterization theorem for positive definite
sequences of the Krawtchouk polynomials}, Australian J. Statist.
\textbf{11} (1969), 29--38.

\bibitem{FeinsilverFitzgerald1996}
P.~Feinsilver and R.~Fitzgerald, \emph{The spectrum of symmetric
Krawtchouk matrices}, Linear Algebra Appl. \textbf{235} (1996), 121--139.
\href{https://doi.org/10.1016/0024-3795(94)00123-5}{doi:10.1016/0024-3795(94)00123-5}.

\bibitem{FeinsilverSFAII2023}
P.~Feinsilver, \emph{Symmetric Functions Algebras (SFA) II: Induced
Matrices}, Journal of Stochastic Analysis \textbf{4} (2023), no.~2,
Article~1.
\href{https://doi.org/10.31390/josa.4.2.01}{doi:10.31390/josa.4.2.01}.

\bibitem{FeinsilverKocik2005}
P.~Feinsilver and J.~Kocik, \emph{Krawtchouk polynomials and Krawtchouk
matrices}, in: \emph{Recent Advances in Applied Probability}
(R.~Baeza-Yates, J.~Glaz, H.~Gzyl, J.~H\"usler, and J.~L.~Palacios, eds.),
Springer, New York, 2005, pp.~115--141.
\href{https://doi.org/10.1007/0-387-23394-6_5}{doi:10.1007/0-387-23394-6\_5}.

\bibitem{GodoyEtAl1998}
E.~Godoy, A.~Ronveaux, A.~Zarzo, and I.~Area, \emph{On the limit
relations between classical continuous and discrete orthogonal
polynomials}, J. Comput. Appl. Math. \textbf{91} (1998), 97--105.
\href{https://doi.org/10.1016/S0377-0427(98)00026-0}{doi:10.1016/S0377-0427(98)00026-0}.

\bibitem{Griffiths1971}
R.~C. Griffiths, \emph{Orthogonal polynomials on the multinomial distribution},
Australian J. Statist. (1971).
\href{https://doi.org/10.1111/j.1467-842X.1971.tb01239.x}{doi:10.1111/j.1467-842X.1971.tb01239.x}.

\bibitem{Griffiths2016}
R.~C. Griffiths, \emph{Multivariate Krawtchouk polynomials and composition
birth and death processes}, Symmetry \textbf{8} (2016), no.~5, article 33.
\href{https://doi.org/10.3390/sym8050033}{doi:10.3390/sym8050033}.

\bibitem{Krawtchouk1929}
M.~Krawtchouk, \emph{Sur une g\'en\'eralisation des polyn\^omes
d'Hermite}, C. R. Acad. Sci. Paris \textbf{189} (1929), 620--622.

\bibitem{MacWilliamsSloane1977}
F.~J. MacWilliams and N.~J.~A. Sloane, \emph{The Theory of
Error-Correcting Codes}, North-Holland, Amsterdam, 1977.

\bibitem{Mehler1866}
F.~G. Mehler, \emph{Ueber die Entwicklung einer Function von beliebig
vielen Variablen nach Laplaceschen Functionen h\"oherer Ordnung},
J. Reine Angew. Math. \textbf{66} (1866), 161--176.
\href{https://doi.org/10.1515/crll.1866.66.161}{doi:10.1515/crll.1866.66.161}.
\end{thebibliography}
\end{document}
